\chapter*{Chapitre 5 — Distribution et émission de chaleur}
\addcontentsline{toc}{chapter}{Chapitre 5 — Distribution et émission de chaleur}

% ============================================================
\section*{Énoncés}
% ============================================================

\exercice{Réseau bitube — sélection des radiateurs d'un appartement}

Un appartement de 4 pièces doit être équipé de radiateurs à eau chaude 70/55\,°C ($T_{amb} = 20$\,°C, exposant $n = 1{,}3$). Les déperditions par pièce sont :

\begin{center}
\begin{tabular}{lcc}
  \toprule
  Pièce & Déperditions (W) & Surface (m²) \\
  \midrule
  Séjour & 1600 & 25 \\
  Chambre 1 & 800 & 12 \\
  Chambre 2 & 700 & 10 \\
  Cuisine & 500 & 8 \\
  \bottomrule
\end{tabular}
\end{center}

\begin{enumerate}
  \item Calculer la puissance nominale nécessaire pour chaque pièce.
  \item Calculer la puissance totale à fournir. Vérifier la cohérence avec le dimensionnement de la chaudière.
\end{enumerate}

% ----------------------------------------------------------

\exercice{Plancher chauffant — dimensionnement complet d'une pièce}

Un salon de 20\,m² présente des déperditions de \SI{900}{\watt}. Revêtement : parquet flottant compatible PCH ($R_{revêt} = \SI{0{,}10}{\square\metre\kelvin\per\watt}$). Le PCH fonctionne en 40/32\,°C, $T_{amb} = 20$\,°C.

\begin{enumerate}
  \item Calculer la puissance surfacique nécessaire.
  \item Calculer le coefficient $B$ et la puissance surfacique disponible avec $p = 150$\,mm.
  \item Calculer la longueur de tube nécessaire.
  \item Vérifier si une seule boucle suffit.
\end{enumerate}

% ----------------------------------------------------------

\exercice{Diagnostic d'un radiateur sous-performant en salle de réunion}

Un technicien mesure sur un radiateur de salle de réunion les valeurs suivantes :
\begin{itemize}
  \item Température entrée eau : 72\,°C
  \item Température sortie eau : 62\,°C
  \item Débit eau mesuré : 60\,L/h
  \item Température ambiante : 20\,°C
\end{itemize}
Le local présente des déperditions de 1500\,W. Le radiateur a une puissance nominale de 2200\,W (en 75/65/20\,°C, $n = 1{,}3$).

\begin{enumerate}
  \item Calculer la puissance thermique réellement délivrée.
  \item Calculer la puissance théorique attendue dans ces conditions (en appliquant la correction d'écart de température logarithmique moyen).
  \item Analyser l'écart entre puissance réelle et théorique, et proposer des pistes de diagnostic.
\end{enumerate}

\newpage
% ============================================================
\section*{Corrigés}
% ============================================================

\subsection*{Corrigé — Exercice 1}

\textit{Question 1 — Écart de température réel :}
\[
  T_{moy} = \frac{70 + 55}{2} = \SI{62{,}5}{\celsius}, \quad \Delta T_{réel} = 62{,}5 - 20 = \SI{42{,}5}{\kelvin}
\]
\[
  \text{Facteur correction} = \left(\frac{42{,}5}{50}\right)^{1{,}3} = (0{,}85)^{1{,}3} = 0{,}826
\]

\begin{methode}[title=Correction de puissance d'un émetteur]
La puissance nominale d'un radiateur est mesurée en conditions normalisées (75/65/20\,°C,
soit $\Delta T_{norm} = 50$\,K). Pour des conditions réelles différentes, on applique :
\[
  \Phi_{réel} = \Phi_{nom} \times \left(\frac{\Delta T_{réel}}{\Delta T_{norm}}\right)^n
\]
Pour dimensionner, on cherche la puissance nominale à sélectionner pour obtenir les déperditions en conditions réelles :
\[
  \Phi_{nom,néc} = \frac{\Phi_{dép}}{\text{facteur correction}}
\]
\end{methode}

\begin{center}
\begin{tabular}{lcc}
  \toprule
  Pièce & Déperditions (W) & $\Phi_{nom,néc}$ = Dép. / 0,826 (W) \\
  \midrule
  Séjour & 1600 & 1937 \\
  Chambre 1 & 800 & 969 \\
  Chambre 2 & 700 & 848 \\
  Cuisine & 500 & 605 \\
  \midrule
  \textbf{Total} & \textbf{3600} & \textbf{4359} \\
  \bottomrule
\end{tabular}
\end{center}

\textit{Question 2 :} La puissance totale de déperditions est 3600\,W. La chaudière devra fournir au minimum 3,6\,kW, majoré des pertes de distribution (5–10\,\%) → chaudière $\geq$ 4\,kW.

\[
  \boxed{P_{chaudière} \geq \SI{4}{\kilo\watt}}
\]

% ----------------------------------------------------------

\subsection*{Corrigé — Exercice 2}

\textit{Question 1 :}
\[
  q_{néc} = \frac{900}{20} = \SI{45}{\watt\per\square\metre}
\]

\[
  \boxed{q_{néc} = \SI{45}{\watt\per\square\metre}}
\]

\textit{Question 2 :}
\[
  B = \frac{1}{0{,}173 + 0{,}10} = \frac{1}{0{,}273} = \SI{3{,}66}{\watt\per\square\metre\per\kelvin}
\]
\[
  T_{moy} = \frac{40 + 32}{2} = \SI{36}{\celsius}, \quad \Delta T = 36 - 20 = \SI{16}{\kelvin}
\]
\[
  q_{disp} = B \times \frac{\Delta T}{p} = 3{,}66 \times \frac{16}{0{,}15} = \SI{390}{\watt\per\square\metre}
\]

La puissance disponible (390\,W/m²) est très largement supérieure à la puissance nécessaire (45\,W/m²). Même avec $p = 300$\,mm, $q = 195$\,W/m² >> 45\,W/m². Choisir $p = \SI{300}{\milli\metre}$.

\textit{Question 3 — Longueur de tube :}
\[
  L_{tube} = \frac{S}{p} = \frac{20}{0{,}30} = \SI{67}{\metre}
\]

\[
  \boxed{L_{tube} = \SI{67}{\metre}}
\]

\textit{Question 4 :} 67\,m < 100\,m → une seule boucle suffit.

% ----------------------------------------------------------

\subsection*{Corrigé — Exercice 3}

\textit{Question 1 — Puissance réelle mesurée :}
\[
  \dot{m} = \frac{60}{3600} \times 1{,}0 = \SI{0{,}0167}{\kilogram\per\second}
\]
\[
  \dot{Q}_{réel} = \dot{m} \times c_p \times \Delta T = 0{,}0167 \times 4186 \times (72 - 62) = \SI{699}{\watt} \approx \SI{700}{\watt}
\]

\[
  \boxed{\dot{Q}_{réel} \approx \SI{700}{\watt}}
\]

\textit{Question 2 — Puissance théorique attendue :}

On applique la correction de puissance par l'exposant $n = 1{,}3$ (conditions normalisées : 75/65/20\,°C, soit $\Delta T_{norm} = 50$\,K) :
\[
  T_{moy} = \frac{72 + 62}{2} = \SI{67}{\celsius}, \quad \Delta T_{réel} = 67 - 20 = \SI{47}{\kelvin}
\]
\[
  \dot{Q}_{théo} = 2200 \times \left(\frac{47}{50}\right)^{1{,}3} = 2200 \times (0{,}94)^{1{,}3}
\]
\[
  (0{,}94)^{1{,}3} = e^{1{,}3 \times \ln(0{,}94)} = e^{1{,}3 \times (-0{,}0619)} = e^{-0{,}0805} \approx 0{,}923
\]
\[
  \dot{Q}_{théo} = 2200 \times 0{,}923 = \SI{2031}{\watt}
\]

\[
  \boxed{\dot{Q}_{théo} \approx \SI{2031}{\watt}}
\]

\textit{Question 3 — Analyse :} Écart important : 700\,W mesuré vs 2031\,W théorique (-65\,\%). Les déperditions de 1500\,W ne sont pas couvertes — le local est en sous-chauffe. Pistes de diagnostic :
\begin{itemize}
  \item Débit très faible (60\,L/h vs débit nominal estimé à 130–180\,L/h) → robinet thermostatique bloqué ou partiellement fermé, ou vanne d'équilibrage surréduite
  \item Présence d'air dans le radiateur → purger le purgeur manuel
  \item Encrassement intérieur (dépôts de boues magnétites) réduisant la section de passage et l'échange thermique
  \item Déséquilibre hydraulique du réseau → vérifier les pressions différentielles avec un manomètre
\end{itemize}
